\documentclass[12pt]{article}
\usepackage[lmargin=1cm, rmargin=2cm,tmargin=2cm,bmargin=2cm]{geometry}
\usepackage{amsmath,amssymb}
\usepackage{enumitem}

\begin{document}

\section*{Hypothesis Testing Question Bank with Answers (2026--27)}

\noindent
These questions are aligned with the Topic 2 inference slides. They keep the technical level light: students should focus on hypotheses, sampling variability, standard errors, p-values, confidence intervals, and interpretation.

\subsection*{1. Naming the hypotheses}

\noindent
\textbf{Question.}
A student newspaper claims that students now spend fewer than 10 hours per week preparing for classes outside the classroom. A sample of students reports an average of 9.1 hours per week.

\begin{enumerate}[label=(\alph*)]
\item What is the null hypothesis?
\item What is the alternative hypothesis if we want to test the newspaper's specific claim?
\item What would the alternative hypothesis be if we only wanted to test whether average study time has changed?
\end{enumerate}

\noindent
\textbf{Answer.}
\begin{enumerate}[label=(\alph*)]
\item $H_0: \mu = 10$ hours per week.
\item $H_A: \mu < 10$ hours per week. This is a one-sided test.
\item $H_A: \mu \ne 10$ hours per week. This is a two-sided test.
\end{enumerate}

\hrule

\subsection*{2. A survey and a one-sided test}

\noindent
\textbf{Question.}
In a survey of $n=400$ residents, $216$ support a proposed rent-control policy. City officials want to know whether support is above $50\%$.

\begin{enumerate}[label=(\alph*)]
\item Write the null and alternative hypotheses.
\item Compute the sample proportion.
\item Using the null standard error
\[
SE_0 = \sqrt{\frac{0.50(1-0.50)}{400}} = 0.025,
\]
compute the test statistic.
\item The one-sided p-value is about $0.055$. What should we conclude at the $5\%$ level?
\end{enumerate}

\noindent
\textbf{Answer.}
\begin{enumerate}[label=(\alph*)]
\item $H_0: p = 0.50$ and $H_A: p > 0.50$.
\item $\hat{p}=216/400=0.54$.
\item
\[
z=\frac{0.54-0.50}{0.025}=1.60.
\]
\item At the $5\%$ level, we do not reject $H_0$ because $0.055>0.05$. The sample leans toward majority support, but the evidence is just short of the usual $5\%$ threshold. This is a good case for reporting the estimate and the p-value, not only the reject/do-not-reject decision.
\end{enumerate}

\hrule

\subsection*{3. Trust in parliament}

\noindent
\textbf{Question.}
Suppose a previous study found that average trust in parliament was $50$ on a 0--100 scale. This year, a survey of $n=100$ students has mean trust $\bar{x}=53$ and standard deviation $s=15$. Use the approximate standard error $SE=s/\sqrt{n}=1.5$.

\begin{enumerate}[label=(\alph*)]
\item Test whether average trust has changed from 50 using a two-sided test.
\item The two-sided p-value for $z=2.00$ is about $0.046$. Interpret it.
\item Give an approximate $95\%$ confidence interval.
\end{enumerate}

\noindent
\textbf{Answer.}
\begin{enumerate}[label=(\alph*)]
\item $H_0:\mu=50$ and $H_A:\mu\ne 50$.
\[
z=\frac{53-50}{1.5}=2.00.
\]
\item If the true mean were 50, a sample mean at least this far from 50 would occur about $4.6\%$ of the time in repeated samples of this size. At the $5\%$ level, this is statistically significant, though only just.
\item
\[
53 \pm 1.96(1.5)=53\pm 2.94,
\]
so the interval is approximately $[50.06,55.94]$.
\end{enumerate}

\hrule

\subsection*{4. Difference in means after a civic-information prompt}

\noindent
\textbf{Question.}
In an experiment, students are randomly assigned to read either a short civic-information prompt or a neutral text before scoring the accuracy of political headlines. Scores run from 0 to 100.

\[
\begin{array}{lccc}
\hline
\text{Group} & n & \text{Mean score} & \text{Standard deviation}\\
\hline
\text{Prompt} & 100 & 68 & 18\\
\text{Control} & 100 & 62 & 20\\
\hline
\end{array}
\]

\begin{enumerate}[label=(\alph*)]
\item What is the estimated difference in means?
\item Use
\[
SE = \sqrt{\frac{18^2}{100}+\frac{20^2}{100}} \approx 2.69
\]
to compute the test statistic for $H_0:\mu_P-\mu_C=0$.
\item The two-sided p-value is about $0.026$. What does this suggest?
\item Why does random assignment matter for the interpretation?
\end{enumerate}

\noindent
\textbf{Answer.}
\begin{enumerate}[label=(\alph*)]
\item The estimated difference is $68-62=6$ points.
\item
\[
z=\frac{6-0}{2.69}\approx 2.23.
\]
\item The result is statistically significant at the $5\%$ level. The data provide evidence that the prompt changed average accuracy scores.
\item Because students were randomly assigned, the prompt and control groups should be comparable before the prompt. That makes the difference in means more credible as evidence of a causal effect of the prompt, not just a pre-existing difference between groups.
\end{enumerate}

\hrule

\subsection*{5. Interpreting a p-value}

\noindent
\textbf{Question.}
A study reports a p-value of $0.03$ for the difference between two campaign messages. Which interpretation is best?

\begin{enumerate}[label=(\Alph*)]
\item There is a $3\%$ probability that the null hypothesis is true.
\item There is a $97\%$ probability that the alternative hypothesis is true.
\item If the null hypothesis were true, results this extreme or more extreme would occur about $3\%$ of the time.
\item The campaign message has a large and important effect.
\end{enumerate}

\noindent
\textbf{Answer.}
Choice (C) is best. A p-value is calculated assuming the null hypothesis is true. It does not give the probability that the null is true, and it does not by itself measure substantive importance.

\hrule

\subsection*{6. Confidence intervals and hypothesis tests}

\noindent
\textbf{Question.}
A policy trial estimates that a tutoring program raises exam scores by $2.8$ points. The approximate $95\%$ confidence interval is $[-1.5,7.1]$.

\begin{enumerate}[label=(\alph*)]
\item Is the result statistically significant at the $5\%$ level for a two-sided test?
\item What does the interval say in plain language?
\item What should we avoid saying?
\end{enumerate}

\noindent
\textbf{Answer.}
\begin{enumerate}[label=(\alph*)]
\item No. The interval includes $0$, so a two-sided test at the $5\%$ level would not reject a zero effect.
\item The best estimate is positive, but the data are compatible with a small negative effect, no effect, or a moderately positive effect.
\item We should avoid saying that the program definitely has no effect. The result is imprecise, not proof of zero.
\end{enumerate}

\hrule

\subsection*{7. Why sample size changes evidence}

\noindent
\textbf{Question.}
Two polls both estimate support for a constitutional reform at $\hat{p}=0.52$. Test $H_0:p=0.50$ against $H_A:p\ne 0.50$.

\begin{enumerate}[label=(\alph*)]
\item Poll A has $n=625$. Using $SE_0=\sqrt{0.25/625}=0.020$, compute the test statistic.
\item Poll B has $n=2500$. Using $SE_0=\sqrt{0.25/2500}=0.010$, compute the test statistic.
\item Explain why the same estimate can produce different evidence.
\end{enumerate}

\noindent
\textbf{Answer.}
\begin{enumerate}[label=(\alph*)]
\item
\[
z=\frac{0.52-0.50}{0.020}=1.00.
\]
This is not statistically significant at the $5\%$ level in a two-sided test.
\item
\[
z=\frac{0.52-0.50}{0.010}=2.00.
\]
This is just significant at about the $5\%$ level in a two-sided test.
\item A larger sample has a smaller standard error. The estimate is the same, but Poll B measures it more precisely, so the same distance from the null looks less compatible with random sampling variation.
\end{enumerate}

\hrule

\subsection*{8. Statistical versus substantive importance}

\noindent
\textbf{Question.}
A very large survey estimates that a new ballot format increases turnout by $0.4$ percentage points, with a p-value below $0.01$.

\begin{enumerate}[label=(\alph*)]
\item Is the effect statistically significant?
\item Does statistical significance automatically mean the effect is substantively large?
\item What should a good answer report?
\end{enumerate}

\noindent
\textbf{Answer.}
\begin{enumerate}[label=(\alph*)]
\item Yes. A p-value below $0.01$ is statistically significant at the $1\%$ level.
\item No. With a very large sample, even a tiny effect can be estimated precisely enough to be statistically significant.
\item A good answer reports the estimate, uncertainty or p-value, the direction of the effect, and whether the size is meaningful in the policy context.
\end{enumerate}

\end{document}
